
\section{Hướng phát triển trong tương lai}
    Mặc dù sản phẩm đã đạt được những kết quả khả quan và đáp ứng đầy đủ các mục tiêu thiết kế cốt lõi của một IoT Gateway dạng module, việc đưa thiết bị vào vận hành thực tế trong môi trường công nghiệp đòi hỏi cần đạt được những tiêu chuẩn khắt khe. Trong tương lai, sản phẩm cần được đánh giá dựa trên các tiêu chuẩn công nghiệp nhằm đạt được những yêu cầu thực tế, mục tiêu có thể đưa sản phẩm ra thị trường.

    % \subsection{Timeline giai đoạn tiếp theo}
    %     \resizebox{\textwidth}{!}{
    %     \begin{tikzpicture}[
    %         node distance = 1.2cm,
    %         line/.style = {ultra thick, black!30, -latex},
    %         dot/.style = {circle, fill=BKblue, minimum size=8pt, inner sep=0pt, outer sep=0pt},
    %         box/.style = {rectangle, draw=BKblue, fill=BKblue!2, rounded corners=3pt, inner sep=10pt, text width=4.5cm, font=\small, align=center},
    %         time/.style = {font=\bfseries\color{BKblue}, align=center, inner sep=0pt}
    %     ]

    %         % Vẽ trục thời gian chính
    %         \draw[line] (-1,0) -- (23,0);

    %         % --- Mốc 1: Hộp TRÊN - Thời gian DƯỚI (Sát trục) ---
    %         \node[dot] (d1) at (0,0) {};
    %         \node[box, above=0.8cm of d1] {Đánh giá prototype, xác định những điểm cần điều chỉnh hoặc tối ưu};
    %         % yshift=-6pt đưa chữ xuống sát mép dưới của chấm tròn
    %         \node[time, anchor=north, yshift=-6pt] at (d1) {Tháng 01/2026};
    %         \draw[black!30] (d1) -- (0,0.8);

    %         % --- Mốc 2: Hộp DƯỚI - Thời gian TRÊN (Sát trục) ---
    %         \node[dot] (d2) at (3,0) {};
    %         \node[box, below=0.8cm of d2] {Nâng cấp, điều chỉnh lại kiến trúc hệ thống};
    %         % yshift=6pt đưa chữ lên sát mép trên của chấm tròn
    %         \node[time, anchor=south, yshift=6pt] at (d2) {Tháng 01/2026};
    %         \draw[black!30] (d2) -- (3,-0.8);

    %         % --- Mốc 3: Hộp TRÊN - Thời gian DƯỚI ---
    %         \node[dot] (d3) at (6,0) {};
    %         \node[box, above=0.8cm of d3] {Thiết kế Kiến trúc hệ thống};
    %         \node[time, anchor=north, yshift=-6pt] at (d3) {Tháng 02/2026};
    %         \draw[black!30] (d3) -- (6,0.8);

    %         % --- Mốc 4: Hộp DƯỚI - Thời gian TRÊN ---
    %         \node[dot] (d4) at (9,0) {};
    %         \node[box, below=0.8cm of d4] {Thiết kế Schematic};
    %         \node[time, anchor=south, yshift=6pt] at (d4) {Tháng 02/2026};
    %         \draw[black!30] (d4) -- (9,-0.8);

    %         % --- Mốc 5: Hộp TRÊN - Thời gian DƯỚI ---
    %         \node[dot] (d5) at (12,0) {};
    %         \node[box, above=0.8cm of d5] {Thiết kế Layout};
    %         \node[time, anchor=north, yshift=-6pt] at (d5) {Tháng 03/2026};
    %         \draw[black!30] (d5) -- (12,0.8);

    %         % --- Mốc 6: Hộp DƯỚI - Thời gian TRÊN ---
    %         \node[dot] (d6) at (15,0) {};
    %         \node[box, below=0.8cm of d6] {Gia công PCB};
    %         \node[time, anchor=south, yshift=6pt] at (d6) {Tháng 03/2026};
    %         \draw[black!30] (d6) -- (15,-0.8);

    %         % --- Mốc 7: Hộp TRÊN - Thời gian DƯỚI ---
    %         \node[dot] (d7) at (18,0) {};
    %         \node[box, above=0.8cm of d7] {Kiểm thử và triển khai};
    %         \node[time, anchor=north, yshift=-6pt] at (d7) {Tháng 04/2026};
    %         \draw[black!30] (d7) -- (18,0.8);

    %         % --- Mốc 8: Hộp DƯỚI - Thời gian TRÊN ---
    %         \node[dot] (d8) at (21,0) {};
    %         \node[box, below=0.8cm of d8] {Báo cáo, phản biện, bảo vệ đồ án};
    %         \node[time, anchor=south, yshift=6pt] at (d8) {Tháng 05/2026};
    %         \draw[black!30] (d8) -- (21,-0.8);

    %     \end{tikzpicture}
    %     }

    % \subsection{Hoàn thiện và khắc phục các hạn chế thiết kế}
    % Dù hệ thống đã vận hành ổn định, quá trình thử nghiệm thực tế cho thấy vẫn tồn tại một số lỗi nhỏ về mặt logic phần cứng và bố trí linh kiện. 
    %     \begin{itemize}
    %         \item \textbf{Tối ưu hóa thiết kế:} Mặc dù đã được thiết kế nhỏ gọn với vỏ hộp và cơ chế lắp đặt module adapter dễ dàng hơn, tuy nhiên việc láp ráp thiết bị vẫn còn khá phức tạp và tốn thời gian. Cần cải tiến thiết kế vỏ hộp và cơ chế lắp đặt để giảm thiểu thời gian lắp ráp, đồng thời tăng tính tiện lợi cho người dùng cuối khi triển khai thực tế.
    %         \item \textbf{Cải thiện Layout PCB:} Tinh chỉnh lại các đường mạch tín hiệu nhạy cảm để giảm thiểu hiện tượng nhiễu xuyên âm (Crosstalk) và cải thiện tính toàn vẹn tín hiệu (Signal Integrity) khi hệ thống hoạt động ở tần số cao.
    %     \end{itemize}

    \subsection{Cải tiến đến tiêu chuẩn công nghiệp}
        Để đưa một thiết bị IoT Gateway ra thị trường, đặc biệt là trong môi trường công nghiệp, sản phẩm phải đáp ứng các tiêu chuẩn khắt khe về an toàn, tương thích điện từ, viễn thông và độ bền cơ học.
        \begin{itemize}
            \item \textbf{Tương thích điện từ (EMC/EMI):} Đảm bảo gateway không phát ra nhiễu điện từ vượt mức cho phép và có khả năng chống chịu nhiễu từ các thiết bị công nghiệp khác (motor, biến tần). (IEC 61000-6-2 và IEC 61000-6-2, \dots)
            \item \textbf{An toàn điện:} Ngăn chặn các rủi ro về giật điện, cháy nổ đối với người vận hành. (IEC/EN 62368-1, \dots)
            \item \textbf{Chuẩn Viễn thông và Vô tuyến:} Đặc biệt quan trọng đối với IoT Gateway tích hợp nhiều giao thức không dây (LoRa, Zigbee, BLE, Wi-Fi) cần đạt tiêu chuẩn kiểm định các dải băng tần. (ETSI EN 300 328 và ETSI EN 300 220, \dots)
            \item \textbf{Giao tiếp công nghiệp có dây:} Đối với các node mạng giao tiếp qua các chuẩn công nghiệp như RS-485, hệ thống cần đáp ứng tiêu chuẩn vật lý (TIA/EIA-485-A) và cách ly (UL 1577).
            \item \textbf{Môi trường và Độ tin cậy:} Đảm bảo thiết bị hoạt động ổn định trong các điều kiện khắc nghiệt của nhà máy, như chuỗi tiêu chuẩn thử nghiệm môi trường gồm nhiệt độ, độ ẩm, rung, sốc cơ học (IEC 60068 Series), tiêu chuẩn chống bụi và nước (IEC 60529) hay tiêu chuẩn về hạn chế các chất độc hại (RoHS).
        \end{itemize}

    % \subsection{Nâng cấp và Mở rộng năng lực Phần cứng}
    % Nhằm đưa sản phẩm từ một bản thử nghiệm thành một thiết bị ở quy mô công nghiệp. Sản phẩm prototype hiện tại chỉ đáp ứng được các chức năng cốt lõi của một gateway, song để sản phẩm thực sự hoàn thiện và có thể sử dụng ngoài thực tế thì vẫn còn nhiều chức năng cần phải được đánh giá và nâng cấp.



    \subsection{Phát triển Firmware chuyên sâu và tối ưu hóa hiệu suất}
    % Hiện tại, phần mềm hệ thống trên cả hai vi điều khiển (Master và Slave) mới chỉ dừng lại ở mức độ nguyên mẫu (Proof of Concept) để chứng minh luồng dữ liệu cơ bản. Để sản phẩm đạt độ hoàn thiện cao và sẵn sàng cho triển khai thực tế, các nội dung sau cần được ưu tiên phát triển:

    % \begin{itemize}
    %     \item \textbf{Hoàn thiện Driver và Stack giao thức cho mạng cảm biến (LAN Domain):}
    %     \begin{itemize}
    %         \item Nâng cấp các giao thức hiện hữu (LoRa, Zigbee, RS485): Hiện tại, các module này mới chỉ được phát triển ở mức độ thử nghiệm (demo) gửi/nhận gói tin đơn giản (raw data). Cần triển khai đầy đủ các tầng của ngăn xếp giao thức (Protocol Stack) như cơ chế địa chỉ hóa (Addressing), quản lý mạng (Network Management), và kiểm tra lỗi gói tin (Packet Error Checking) để đảm bảo độ tin cậy.
    %         \item Phát triển mới các giao thức chưa được hỗ trợ: Tiến hành lập trình driver và tích hợp các giao thức đã có trong thiết kế nhưng chưa được hiện thực hóa, bao gồm Z-Wave, Thread và Bluetooth Low Energy (BLE). Việc này đòi hỏi tích hợp các thư viện stack chuẩn và cấu hình tài nguyên phần cứng (Timer, Interrupt) phù hợp trên LAN-MCU.
    %     \end{itemize}
        
    %     \item \textbf{Mở rộng khả năng kết nối mạng và giao thức ứng dụng (WAN Domain):}
    %     \begin{itemize}
    %         \item Đa dạng hóa kênh kết nối Internet: Hiện tại hệ thống chưa hỗ trợ Ethernet và NB-IoT. Cần phát triển driver cho module Ethernet (như W5500 hoặc EMAC tích hợp) và module NB-IoT để tăng tính dự phòng đường truyền khi Wi-Fi hoặc LTE gặp sự cố.
    %         \item Bổ sung giao thức tầng ứng dụng: Ngoài MQTT hiện đang hoạt động, cần phát triển thêm các client cho HTTP/HTTPS (phục vụ gọi REST API) và CoAP (Constrained Application Protocol) để tương thích với các nền tảng IoT khác nhau và tối ưu hóa băng thông cho các mạng di động.
    %     \end{itemize}
        
    %     \item \textbf{Cải tiến kiến trúc luồng dữ liệu và độ tin cậy hệ thống:}
    %     \begin{itemize}
    %         \item Tối ưu hóa luồng Gửi/Nhận (TX/RX Pipeline): Thiết kế lại cơ chế đệm (Buffering) và quản lý hàng đợi (Queue management) giữa các tác vụ (Tasks) trong FreeRTOS. Chuyển đổi từ cơ chế xử lý tuần tự hoặc blocking sang cơ chế hướng sự kiện (Event-driven) hoặc sử dụng DMA (Direct Memory Access) để tăng thông lượng và giảm tải cho CPU.
    %         \item Xây dựng cơ chế xử lý lỗi (Error Handling) toàn diện: Hiện tại cơ chế xử lý lỗi còn sơ sài. Cần bổ sung các trình phục vụ ngắt và giám sát trạng thái (State Machine) để tự động phát hiện và khôi phục (Self-healing) khi xảy ra các lỗi như: mất kết nối module, treo bus SPI/I2C, hoặc tràn bộ nhớ đệm.
    %     \end{itemize}
    % \end{itemize}

    Dựa trên nền tảng vận hành ổn định và kiến trúc hoàn thiện của Phiên bản 1.0.0, các định hướng phát triển tiếp theo tập trung vào việc tối ưu hóa tài nguyên và nâng cao năng lực xử lý để đáp ứng các kịch bản công nghiệp quy mô lớn với mật độ dữ liệu cao:

    \begin{itemize}
        \item \textbf{Nâng cấp toàn diện các ngăn xếp giao thức (Protocol Stacks):} Mặc dù hệ thống đã hỗ trợ đa dạng giao thức Sensing (Zigbee, LoRa, RS-485), việc triển khai hiện tại chủ yếu tập trung vào lớp vận chuyển dữ liệu thô. Để sẵn sàng cho các hệ thống Enterprise, cần tích hợp đầy đủ các tầng quản lý mạng chuyên sâu bao gồm: cơ chế định danh thiết bị quy mô lớn, quản lý cấu trúc Mesh phức tạp và các thuật toán kiểm tra lỗi gói tin nâng cao để duy trì tính toàn vẹn dữ liệu tuyệt đối.
        
        \item \textbf{Tối ưu hóa pipeline dữ liệu bằng kiến trúc hướng sự kiện (Event-driven):} Kết quả đo lường hiệu năng tại Mục~\ref{sec:benchmark} cho thấy kênh SPI inter-MCU đạt thông lượng thực tế $\approx 21\,\text{Mbps}$, dư hơn 190 lần so với tổng nhu cầu $109{,}3\,\text{kbps}$ của ba giao thức vô tuyến đồng thời, do đó bus SPI không phải điểm nghẽn. Tuy nhiên, rủi ro tràn bộ đệm vẫn tồn tại khi LAN-MCU phải phục vụ các ngắt ngoại vi liên tục ở tần suất cao. Do đó, việc chuyển đổi sang kiến trúc hướng sự kiện kết hợp với cơ chế truy cập bộ nhớ trực tiếp (DMA) là lựa chọn kỹ thuật tất yếu để giải phóng tài nguyên CPU, đảm bảo duy trì tỷ lệ truyền gói tin thành công (PDR) $>$ 85\% ngay cả dưới tải cực nặng.
        
        \item \textbf{Mở rộng các giao thức ứng dụng và bảo mật kết nối:} Để tăng cường khả năng tương thích với các nền tảng Cloud công nghiệp khác nhau, hệ thống cần tích hợp thêm các Client chuyên dụng cho HTTP/HTTPS và CoAP bên cạnh MQTT hiện có. Song song đó, cơ chế \textit{Safety FOTA} cần được nâng cấp để hỗ trợ cập nhật song song qua Ethernet, giúp tối ưu hóa băng thông và giảm thiểu thời gian dừng hệ thống (downtime) so với phương thức Wi-Fi NAT hiện tại.
        
        \item \textbf{Xây dựng cơ chế tự phục hồi:} Hệ thống cần bổ sung các máy trạng thái giám sát chuyên sâu để tự động phát hiện và khôi phục các lỗi logic ngoại vi. Việc tự động hóa quy trình xử lý lỗi này sẽ đảm bảo tính sẵn sàng cao của thiết bị mà không cần sự can thiệp vật lý tại hiện trường.
    \end{itemize}

    \subsection{Tối ưu hóa chi phí và Hướng tới sản xuất}
    % Tối ưu hóa chi phí thiết kế phần cứng bằng cách phân tích và lựa chọn các linh kiện tương đương có giá thành tốt hơn hoặc các thiết kế tương đương nhưng vẫn đảm kỹ thuật và tiêu chuẩn công nghiệp, đánh giá về chi phí gia công.
    Quy trình tối ưu hóa chi phí thiết kế phần cứng được thực hiện thông qua việc phân tích và lựa chọn hệ thống linh kiện thay thế có thông số kỹ thuật tương đương, đảm bảo duy trì độ tin cậy theo tiêu chuẩn công nghiệp trong khi vẫn giảm thiểu tổng giá thành linh kiện (BOM). Song song đó, việc thu nhỏ kích thước PCB được tính toán kỹ lưỡng để trực tiếp cắt giảm chi phí phôi mạch in và chi phí gia công lắp ráp hàng loạt, tạo lợi thế về kinh tế cho sản phẩm mà không làm suy giảm hiệu năng vận hành hay tính toàn vẹn của tín hiệu.

